\section{Product projections and localization}
\label{sec:localization}

The parity proof will use localization for a single output projection.
We prove the stronger uniform statement for every product projection,
because its induction exchanges the projected and counted registers.

All operators act on \(\mathcal H_N=(\mathbb C^2)^{\otimes N}\), and
\(\norm{\cdot}\) denotes the operator norm. Identities on unmentioned
qubits are implicit. Spectral projections are written in brackets;
for example, \([M\ge r]\) is the projection onto the sum of eigenspaces
of a self-adjoint operator \(M\) with eigenvalues at least \(r\).
All logarithms are natural.

A \emph{count} is an operator
\begin{equation}
\label{eq:count}
 M=\sum_{i\in S}\bigl(I-\proj{v_i}\bigr)_i,
 \qquad S\subseteq\{1,\ldots,N\},\quad \norm{v_i}=1.
\end{equation}
The subscript indicates the qubit on which the summand acts.
The reference vector \(v_i\in\mathbb C^2\) can be chosen separately
at each qubit. The summands commute, and the spectrum of \(M\) is
contained in \(\{0,\ldots,|S|\}\).
A \emph{product projection} is a tensor product of rank-one one-qubit
projections on a subset, with identity on the complement. Its support
is the selected subset. The empty
product is \(I\). These are precisely the projections \([M=0]\).
In particular, a product projection need not have rank one on
\(\mathcal H_N\).

We use a slightly enlarged gate set of product-state reflections, as in
\cite[Section~2 and Proposition~2.1]{Rosenthal2021}.
Let \(\cU_d(N)\) be the set of
unitaries
\begin{equation}
\label{eq:normal-form}
 U=(L_dR_d)\cdots(L_1R_1)L_0,
\end{equation}
where each \(L_j\) is a product of one-qubit unitaries and each \(R_j\)
is a layer of gates \(I-2A\) with disjoint supports, with \(A\) a product
projection on the gate's qubits. Identity layers are allowed.
The class \(\cU_0(N)\) consists of products of one-qubit unitaries.

This model contains the circuits of \Cref{thm:parity}. Indeed, if
\(C\) projects onto all controls being \(1\), a Toffoli gate equals
\[
 (I-C)\otimes I+C\otimes X
 =I-2C\otimes\proj{-},
 \qquad \ket-=\frac{\ket0-\ket1}{\sqrt2}.
\]
Grouping the local gates in each parallel step gives
\eqref{eq:normal-form} with a number of reflection layers bounded by
the original depth. Discarded qubits can be retained as idle wires
without changing any later measurement probability. Here $d$ counts
reflection layers in the enlarged model; it is at most the number of
physical layers in the original circuit. The arbitrary local layers in
\eqref{eq:normal-form} are included in this definition, rather than
charged separately to $d$.

We will also use closure under inversion at the same depth. Each
\(R_j\) is self-adjoint, since it is a product of commuting
self-adjoint reflections. Thus
\begin{equation}
\label{eq:inverse}
 U^\dagger=L_0^\dagger R_1L_1^\dagger\cdots R_dL_d^\dagger
 \in\cU_d(N).
\end{equation}
This is the same normal form with the layer indices reversed: the new
local layers are $L_d^\dagger,\ldots,L_0^\dagger$, in temporal order,
and the new reflection layers are $R_d,\ldots,R_1$. No factors on
overlapping supports have been commuted.

\begin{theorem}[Localization of conjugated product projections]
\label{thm:localization}
Fix a nonnegative integer \(d\) and real numbers \(0\le s<t\).
For all sufficiently large \(N\), every \(U\in\cU_d(N)\) and every
pair of counts \(M,D\) on \(\mathcal H_N\) satisfy
\begin{equation}
\label{eq:localization}
 \norm{[M\ge N^t]\,U[D=0]U^\dagger[M=0]}
 \le \exp(-N^s).
\end{equation}
The sufficient lower bound on \(N\) depends only on \(d,s,t\), not
on the circuit, the counted subsets, or their reference vectors.
\end{theorem}

Write \(\cB_d(s,t)\) for the assertion of
\Cref{thm:localization} at the specified \(d,s,t\), including its
uniformity. If \(t>1\), it is immediate for \(N>1\), since
\([M\ge N^t]=0\). We prove the theorem by induction on \(d\), assuming
all fixed exponent pairs at the preceding depth.
The counts \(M,D\) are allowed to be different because the proof
will exchange their roles when applying \eqref{eq:inverse}.

The next two sections supply the approximation and depth-reduction
estimates. The proof of \Cref{thm:localization} is completed in
\Cref{sec:bootstrap}.
